\documentclass[10pt,a4paper]{article}
\usepackage[utf8]{inputenc}
\usepackage[vietnamese]{babel}
\usepackage{amsmath,amssymb,amsfonts}
\usepackage{geometry}
\geometry{a4paper, left=10mm, right=10mm, top=10mm, bottom=10mm}
\usepackage{tikz}
\usetikzlibrary{shapes,arrows.meta,positioning,calc}
\usepackage[most]{tcolorbox}
\usepackage{array,booktabs,multicol}
\usepackage{fontawesome}
\usepackage{helvet}
\renewcommand{\familydefault}{\sfdefault}
\usepackage{parskip}

\definecolor{primary}{RGB}{112, 48, 160}
\definecolor{teal}{RGB}{0, 128, 128}
\definecolor{orange}{RGB}{230, 108, 44}
\definecolor{green}{RGB}{46, 139, 87}
\definecolor{rose}{RGB}{192, 0, 0}

\pagestyle{empty}

\begin{document}

% --- HEADER BANNER ---
\begin{tcolorbox}[
  enhanced,
  colback=violet!8!white,
  colframe=violet!80!black,
  arc=2.5mm, boxrule=1pt,
  left=3mm, right=3mm, top=2mm, bottom=2mm
]
  \begin{tabular*}{\textwidth}{@{\extracolsep{\fill}} lr}
    {\small\bfseries\color{violet!85!black}\faBook\ SỔ TAY KIẾN THỨC TOÁN 2} & 
    {\small\bfseries\color{teal!85!black}\faGraduationCap\ CHƯƠNG TRÌNH GDPT 2018}
  \end{tabular*}
  \vspace{-1mm}
  \begin{center}
    \rule{0.98\textwidth}{0.6pt}
  \end{center}
  \vspace{-2mm}
  \centering
  {\LARGE\bfseries\color{violet!90!black}\faBookmark\ BÀI 70: ÔN TẬP PHÉP CỘNG, PHÉP TRỪ TRONG PHẠM VI 1000}\\[1mm]
  {\itshape\footnotesize\color{black!75} Tổng hợp kĩ năng cộng trừ có nhớ/không nhớ trong phạm vi 1000}
\end{tcolorbox}

\vspace{-2mm}

\noindent
\begin{minipage}[t]{0.49\textwidth}\vspace{0pt}
\begin{tcolorbox}[
  enhanced,
  colback=violet!4!white,
  colframe=violet!75!black,
  arc=2mm, boxrule=0.8pt,
  left=2mm, right=2mm, top=1.8mm, bottom=1.8mm,
  title={\footnotesize\bfseries\color{white}\faKey\ 1. Bảng quy tắc tính toán chuẩn mực}
]
{\footnotesize
\begin{itemize}\setlength{\itemsep}{1pt}\setlength{\leftskip}{-2.5mm}
  \item \textbf{Quy tắc 1: Đặt tính thẳng hàng}\\
  Chữ số hàng đơn vị, chục, trăm phải thẳng cột tuyệt đối.
  \item \textbf{Quy tắc 2: Thứ tự tính toán}\\
  Luôn tính từ \textbf{phải sang trái}:
  $$\text{Đơn vị} \longrightarrow \text{Chục} \longrightarrow \text{Trăm}$$
  \item \textbf{Quy tắc 3: Xử lí có nhớ}\\
  - Cộng có nhớ: Thêm $1$ vào hàng liền kề bên trái.\\
  - Trừ có nhớ: Thêm $1$ vào số trừ ở hàng liền kề bên trái.
\end{itemize}
}
\end{tcolorbox}
\end{minipage}
\hfill
\begin{minipage}[t]{0.49\textwidth}\vspace{0pt}
\begin{tcolorbox}[
  enhanced,
  colback=teal!4!white,
  colframe=teal!75!black,
  arc=2mm, boxrule=0.8pt,
  left=2mm, right=2mm, top=1.8mm, bottom=1.8mm,
  title={\footnotesize\bfseries\color{white}\faSliders\ 2. Các dạng toán trọng tâm học kì 2}
]
{\footnotesize
\begin{itemize}\setlength{\itemsep}{1pt}\setlength{\leftskip}{-2.5mm}
  \item \textbf{Dạng 1: Tính nhẩm số tròn trăm, tròn chục:}\\
  $500 + 400 = 900$; \quad $1\,000 - 300 = 700$.
  \item \textbf{Dạng 2: Tính giá trị của biểu thức:}\\
  Thực hiện từ trái qua phải: $250 + 130 - 80 = 300$.
  \item \textbf{Dạng 3: Điền chữ số còn thiếu vào phép tính.}\\
  \item \textbf{Dạng 4: Giải toán có lời văn:}\\
  Toán về nhiều hơn, ít hơn hoặc bài toán hai bước tính.
\end{itemize}
}
\end{tcolorbox}
\end{minipage}

\vspace{-1.5mm}
\begin{tcolorbox}[
  enhanced,
  colback=orange!4!white,
  colframe=orange!80!black,
  arc=2mm, boxrule=0.8pt,
  left=2mm, right=2mm, top=1.5mm, bottom=1.5mm,
  title={\footnotesize\bfseries\color{white}\faSitemap\ 3. Sơ đồ các bẫy thường gặp và bí quyết xử lí nhanh}
]
\centering
\begin{tikzpicture}[x=1.65cm, y=0.75cm]
  \tikzset{arr/.style={-{Stealth[scale=0.8]}}, every node/.append style={align=center}}
  \node[rectangle, rounded corners=2pt, draw=rose!80!black, fill=rose!15, font=\scriptsize\bfseries] (b) at (0,0) {CÁC BẪY HAY GẶP\\1. Quên số nhớ\\2. Đặt lệch cột\\3. Nhầm số trừ};
  \node[rectangle, rounded corners=2pt, draw=green!70!black, fill=green!15, font=\scriptsize\bfseries] (g) at (3.5,0) {BÍ QUYẾT XỬ LÍ\\1. Chấm điểm nhớ\\2. Dóng thẳng cột\\3. Thử lại bằng phép tính ngược};

  \draw[arr, very thick, red!85!black] (b) -- node[above, font=\tiny\bfseries] {Khắc phục bằng} (g);
\end{tikzpicture}
\end{tcolorbox}

\vspace{-1.5mm}
\noindent
\begin{minipage}[t]{0.49\textwidth}\vspace{0pt}
\begin{tcolorbox}[
  enhanced,
  colback=green!4!white,
  colframe=green!60!black,
  arc=2mm, boxrule=0.8pt,
  left=2mm, right=2mm, top=1.8mm, bottom=1.8mm,
  title={\footnotesize\bfseries\color{white}\faCheckCircle\ 4. Bài tập thực hành mẫu}
]
{\footnotesize
\textbf{1. Đặt tính rồi tính:}\\
- $364 + 273 = \mathbf{637}$ (nhớ sang hàng trăm: $6+7=13$).\\
- $825 - 341 = \mathbf{484}$ (nhớ sang hàng trăm: $12-4=8$).\\
\textbf{2. Bài toán có lời văn:}\\
Trường Tiểu học có $450$ học sinh nam, số học sinh nữ ít hơn học sinh nam $35$ bạn. Hỏi trường có bao nhiêu học sinh nữ?\\
\textit{Giải:} Số học sinh nữ của trường là:\\
$450 - 35 = \mathbf{415}$ (học sinh).
}
\end{tcolorbox}
\end{minipage}
\hfill
\begin{minipage}[t]{0.49\textwidth}\vspace{0pt}
\begin{tcolorbox}[
  enhanced,
  colback=rose!4!white,
  colframe=rose!70!black,
  arc=2mm, boxrule=0.8pt,
  left=2mm, right=2mm, top=1.8mm, bottom=1.8mm,
  title={\footnotesize\bfseries\color{white}\faExclamationTriangle\ 5. Lời dặn dò thi học kì}
]
{\footnotesize
\begin{itemize}\setlength{\itemsep}{1.5pt}\setlength{\leftskip}{-2.5mm}
  \item Khi làm bài thi, luôn dùng nháp tính cẩn thận trước khi viết vào bài làm.
  \item Viết chữ số rõ ràng, ngay ngắn, không viết số $0$ giống số $6$, số $1$ giống số $7$.
  \item Đọc kĩ câu hỏi xem đề bài hỏi học sinh nữ hay hỏi \textbf{"tất cả học sinh toàn trường"}.
\end{itemize}
}
\end{tcolorbox}
\end{minipage}


\end{document}
